\section{Introduction}

Point clouds, i.e. sets of points in some \(D\)-dimensional space, are a
versatile model for physical objects in many disciplines, from laser
scans in computer vision to particle models of molecules in structural
biology. Typical operations on pairs of point clouds, such as their
union, intersection, or more complex comparisons, rely on a common
reference frame. However, the observation process generating the point
coordinates often cannot easily provide it. Hence, \emph{registration}
of point clouds (also known as \emph{alignment} or \emph{superimposing})
is a key tool at the beginning of many analysis workflows. This can
consist of either \emph{rigid} registration (one point cloud can be
rotated and translated to match the other) or \emph{non-rigid}
registration (more intricate transformations are allowed).

For performance critical computations involving point clouds,
researchers and practitioners might turn to the fast and
research-appropriate language Julia. While different packages in its
ecosystem provide implementations of basic registration algorithms,
there is no comprehensive library covering the diverse use cases of
rigid and non-rigid registration, to our knowledge. The presented
package aims to fill this gap.

\section{Scope of the package}

PointCloudRegistration.jl addresses applications within structural
biology and beyond concerning point cloud data. Point clouds
representing biological macromolecules were a prime consideration during
the development but no biological features are specifically assumed.
This also means that no surface or normal information is required or
used.

The package is intended for the following general setting: Given two
weighted \(D\)-dimensional point clouds, a \emph{target} \(X\) with
points \({\bm{x}}_{1},\ldots,{\bm{x}}_{I} \in {\mathbb{R}}^{D}\)
and weights \(p_{1},\ldots,p_{I}\) as well as a \emph{source} with
points \({\bm{y}}_{1},\ldots,{\bm{y}}_{J} \in {\mathbb{R}}^{D}\)
and weights \(q_{1},\ldots,q_{J}\), we want to find a rigid or non-rigid
transformation that maps the source onto the target in some optimal way.

In summary, these features are provided by the package: Rigid
registration with known point-to-point correspondences via minimizing
the \emph{root mean square displacement}, the \emph{mean absolute
displacement}, and the \emph{Geman-McClure loss}. Rigid registration
with unknown correspondences via maximizing the \emph{kernel
correlation} and the \emph{iterative closest point} method. Non-rigid
registration via \emph{Bayesian coherent point drift}, \emph{divergence
free shape interpolation}, \emph{optimal transport}, and \emph{kernel
correlation maximization with a Gaussian network model}. Thinning of
point clouds via \emph{\(k\)-means} and \emph{DP-means}.

\section{Rigid registration}
